% SPDX-License-Identifier: GPL-3.0-or-later
% Copyright (C) 2026 Alessandro Greco
\chapter{Анонимность, наблюдение и метаданные}\label{cap:01-anonimato}
\aggiornatoal{23 сентября 2026 г.}{данные Tor Metrics по 20 сентября 2026 г.}

\begin{obiettivi}
К концу этой главы вы сможете:
\begin{itemize}[nosep]
  \item различать конфиденциальность, анонимность, псевдонимность, несвязываемость и ненаблюдаемость;
  \item объяснять, почему метаданные коммуникации способны раскрывать не меньше, а иногда и больше, чем её содержание;
  \item описывать анализ трафика и объяснять, почему одного шифрования недостаточно для защиты от него;
  \item определять, кто использует Tor и почему разнообразие пользователей само по себе является формой защиты.
\end{itemize}
\end{obiettivi}

Эта книга посвящена системе, предназначенной для защиты особого вида информации: не того, \emph{что} мы говорим в сети, а того, \emph{с кем} мы общаемся. Прежде чем разбирать устройство системы, необходимо понять решаемую ею задачу и то, почему она отличается от более привычной задачи, которую решает шифрование.

% ---------------------------------------------------------------------------
\section{Содержание и метаданные}\label{sec:anon-metadati}

Когда два компьютера обмениваются данными через Интернет, каждый пакет несёт два вида информации. \textbf{Содержание} --- это то, чем обмениваются приложения: текст сообщения или содержимое веб-страницы. \textbf{Метаданные} описывают \emph{сам факт коммуникации}: кто отправляет данные, кому, когда, как часто и в каком объёме. Часть метаданных необходима для работы сети, поскольку маршрутизаторы должны знать, куда доставлять пакеты. IP-адрес назначения и IP-адрес источника, на который следует отправлять ответы, передаются открытым текстом в каждом пакете.

Шифрование, ставшее повсеместным благодаря HTTPS, защищает содержание, но оставляет метаданные видимыми. Любой, кто наблюдает сеть, --- от интернет-провайдера пользователя до оператора крупной точки обмена трафиком --- всё ещё может видеть, что пользователь подключается к определённому сайту, когда устанавливается соединение, сколько оно длится и какой объём трафика создаёт.

\begin{figure}[htbp]
\centering
% SPDX-License-Identifier: GPL-3.0-or-later
% Copyright (C) 2026 Alessandro Greco
%% Русское альтернативное описание EPUB.
\altfig{Схема HTTPS-соединения: содержимое показано в блоках с пометкой о шифровании, а окружающие метаданные остаются видимыми открытым текстом, включая IP-адреса источника и назначения, временные характеристики, размеры и частоту пакетов.}
% Что защищает HTTPS и что всё ещё видит наблюдатель.
\begin{tikzpicture}[font=\sffamily\scriptsize,
  riga/.style={draw,minimum height=0.65cm,align=center},
  visibile/.style={riga,fill=avviso!12,draw=avviso!60},
  protetto/.style={riga,fill=torverde!15,draw=torverde!60!black}]
  \node[client,minimum width=1.6cm] (u) at (0,0) {клиент};
  \node[server,minimum width=1.6cm] (s) at (10.6,0) {сайт};
  \draw[very thick,torgrigio!50] (u) -- (s);
  \node[font=\sffamily\scriptsize,text=torgrigio] at (5.3,0.45) {HTTPS};
  \node[visibile,text width=1.9cm] (a) at (2.5,-1.5) {IP-адрес\\отправителя};
  \node[visibile,text width=1.9cm] (b) at (4.6,-1.5) {IP-адрес\\сайта};
  \node[visibile,text width=1.9cm] (c) at (6.7,-1.5) {время и\\длительность};
  \node[visibile,text width=1.9cm] (d) at (8.8,-1.5) {переданные\\байты};
  \node[protetto,text width=4.2cm] (e) at (3.4,-2.9) {запрошенная страница и содержимое};
  \node[protetto,text width=4.2cm] (f) at (7.9,-2.9) {учётные данные, текст сообщений};
  \node[text=avviso] at (0.5,-1.5) {\textbf{метаданные}};
  \node[text=avviso,font=\sffamily\tiny] at (0.5,-1.85) {открытый текст};
  \node[text=torverde!45!black] at (0.5,-2.9) {\textbf{содержимое}};
  \node[text=torverde!45!black,font=\sffamily\tiny] at (0.5,-3.25) {зашифровано};
  \foreach \n in {a,b,c,d} \draw[dotted,avviso!70] (\n.north) -- (\n.north |- u.south);
\end{tikzpicture}

\caption{Что может видеть наблюдатель при HTTPS-соединении. Содержание зашифровано, но метаданные остаются видимыми; именно эту информацию и пытается защищать Tor.}
\label{fig:metadati}
\end{figure}

\subsection{Что раскрывают метаданные}

Метаданные легко принять за безобидную информацию: \enquote{Они знают только, с кем я разговариваю, но не знают, что я говорю}. Это впечатление обманчиво. Достаточно нескольких примеров:
\begin{itemize}
  \item человек в два часа ночи посещает сайт центра помощи жертвам домашнего насилия, а затем сайт адвоката по бракоразводным делам;
  \item за несколько дней до публикации расследования журналист обменивается сообщениями с сотрудником министерства, которого касается расследование;
  \item сотрудник регулярно посещает сайт профсоюза или сайт ассоциации пациентов, посвящённой определённому заболеванию.
\end{itemize}
Ни в одном из этих случаев не требуется читать содержание, чтобы узнать очень многое. У метаданных есть ещё одно преимущество для тех, кто их собирает: они структурированы, их легко хранить и анализировать в больших масштабах, а также сопоставлять с другими данными. Генерал Майкл Хейден, бывший директор как NSA, так и CIA, выразил эту мысль в ставшей известной реплике во время дискуссии в Университете Джонса Хопкинса 1 апреля 2014 года: \enquote{Мы убиваем людей на основании метаданных} (перевод), после чего сразу уточнил: \enquote{но это не то, что мы делаем с этими метаданными}~\cite{hayden2014metadata}. Уточнение не меняет сути: те, кто собирает метаданные, считают их достаточными для решений такой степени серьёзности.

% ---------------------------------------------------------------------------
\section{Анализ трафика}\label{sec:anon-traffico}

\textbf{Анализ трафика} изучает метаданные коммуникаций, чтобы делать выводы, не читая само содержание. Он так же стар, как телекоммуникации. Даже во время войн XX века наблюдение за тем, как часто передаёт вражеская радиостанция, с кем она связывается и в каком темпе работает, позволяло выявлять структуру армии, не расшифровывая сообщения.

В Интернете анализ трафика имеет две основные формы:
\begin{description}
  \item[Анализ адресов] Наблюдатель на линии связи видит IP-адреса источника и назначения. Это простейшая форма и именно с ней Tor борется наиболее непосредственно: ни одна отдельная точка сети не видит обе конечные стороны соединения.
  \item[Анализ объёма и времени] Даже не имея полной адресной информации, наблюдатель может сравнивать потоки по объёму и временным характеристикам. Если пользователь загружает 3\,MB в 10:15:02, а сервер в другой части сети в 10:15:02 отправляет 3\,MB, весьма вероятно, что наблюдения относятся к одной и той же коммуникации. Такая \textbf{корреляция трафика} является главным ограничением Tor, и мы будем неоднократно к ней возвращаться (\cref{cap:04-modello-minaccia,cap:26-attacchi}).
\end{description}

% ---------------------------------------------------------------------------
\section{Точная терминология}\label{sec:anon-termini}

Слово \enquote{анонимность} используется в нескольких значениях, и смешение этих значений неоднократно приводило к ошибкам в проектировании подобных систем. Научная литература, в частности терминология, предложенная Пфитцманном и Хансеном~\cite{pfitzmann2010terminology}, различает несколько понятий:
\begin{description}
  \item[Конфиденциальность] Неавторизованная сторона не может прочитать содержание. Именно это обеспечивает шифрование.
  \item[Анонимность] Субъект не может быть идентифицирован внутри группы субъектов, называемой \textbf{множеством анонимности}. Анонимность никогда не бывает абсолютной: она существует \enquote{среди} других людей. Чем больше и разнообразнее множество, тем выше защита.
  \item[Псевдонимность] Человек действует под идентификатором, отличным от собственного имени, причём этот идентификатор может быть постоянным. Учётная запись с вымышленным именем --- это псевдоним, а не анонимность: все её действия можно связать между собой.
  \item[Несвязываемость] Наблюдатель не может определить, связаны ли два интересующих объекта --- например, два сообщения, пользователь и веб-сайт или два посещения. Именно это свойство Tor пытается обеспечить между пользователями и их пунктами назначения.
  \item[Ненаблюдаемость] Наблюдатель не может даже определить, \emph{происходит ли} коммуникация. Это значительно более сильное свойство, которого Tor не обеспечивает: наблюдатель, следящий за пользователем, может увидеть использование Tor, если только не применяется подключаемый транспорт (\cref{cap:16-censura}).
\end{description}

\begin{storia}[Отделение идентификации от маршрутизации]
Трое исследователей Naval Research Laboratory, придумавших onion routing в 1995 году, --- Дэвид Голдшлаг, Майкл Рид и Пол Сайверсон --- предпочитали не говорить об анонимности. Их цель, как они писали в статье 1996 года, к которой Сайверсон позднее вернулся в воспоминаниях об этом периоде, состояла в том, чтобы \enquote{отделить идентификацию от маршрутизации}~\cite{syverson2011peel}. Пользователю может понадобиться подтвердить свою личность системе, с которой он взаимодействует, например рабочему компьютеру, не раскрывая каждой промежуточной точке сети, что кто-то из его местоположения подключается именно к этому компьютеру.
\end{storia}

% ---------------------------------------------------------------------------
\section{Кому нужна анонимность?}\label{sec:anon-utenti}

По оценкам Tor Metrics, в августе и сентябре 2026 года к сети Tor ежедневно подключались от 2,8 до 3,5 млн пользователей без использования мостов, в среднем около 3,2 млн. Для Италии оценка находилась в диапазоне от \num{44000} до \num{83000} пользователей в день, в среднем около \num{60000}. Эти оценки характеризуются значительной неопределённостью, поскольку сеть оценивает число пользователей, а не подсчитывает их напрямую (\cref{cap:24-metrics}). Кто эти люди? Они принадлежат к очень разным группам:
\begin{itemize}
  \item обычные пользователи, которые не хотят, чтобы рекламные компании, интернет-провайдеры или операторы используемых ими сетей отслеживали их;
  \item журналисты и их источники, активисты и правозащитники, особенно в странах с авторитарными режимами;
  \item люди, живущие в условиях цензуры и использующие Tor для доступа к заблокированным сайтам (\cref{cap:16-censura});
  \item жертвы насилия или преследования, которым необходимо скрывать своё местоположение;
  \item компании, не желающие раскрывать конкурентам направления своих исследований;
  \item сотрудники правоохранительных органов и военнослужащие, которым требуется посещать сайты или собирать информацию, не раскрывая свою личность: onion routing возник в лаборатории ВМС США именно для этой цели.
\end{itemize}
Небольшая доля пользователей применяет Tor для незаконной деятельности, как это происходит с любым средством связи. Публичная дискуссия обычно сосредоточена именно на этой группе. В части~IV и \cref{cap:25-legale} мы рассмотрим эту тему на основе фактических данных и конкретных случаев.

\subsection{Анонимность любит компанию}

Разнообразие пользователей важно и по более глубокой причине. Система анонимности, которой пользуется только одна группа --- например, военные, --- не защищала бы никого в этой группе: наблюдатель, увидев подключение к сети, сразу понял бы, что перед ним военный, а любой наблюдатель исходящего трафика знал бы, что тот пришёл от военных. Сайверсон приводит именно этот аргумент применительно к сети, которую ВМС США могли бы построить исключительно для себя: всё, что выходило бы из неё, распознавалось бы как трафик ВМС~\cite{syverson2011peel}. Анонимность каждого человека зависит от количества других пользователей и от их разнообразия. Эту идею кратко выражает фраза \enquote{анонимность любит компанию} (в оригинале \emph{Anonymity Loves Company}), которая также стала названием статьи Динглдайна и Мэтьюсона~\cite{dingledine2006company}. Именно поэтому Tor задуман как бесплатная и открытая система, доступная каждому: чем больше людей ею пользуются, тем лучше она защищает всех.

% ---------------------------------------------------------------------------
\section*{Итоги}
\begin{itemize}
  \item Шифрование защищает содержание; метаданные --- кто, с кем, когда и в каком объёме общается --- остаются видимыми и могут раскрывать очень многое.
  \item Анализ трафика извлекает сведения из адресов, объёмов и временных характеристик коммуникации.
  \item Анонимность означает невозможность идентификации внутри множества; её необходимо отличать от псевдонимности, несвязываемости и ненаблюдаемости.
  \item Tor обслуживает миллионы очень разных пользователей, и это разнообразие является частью его защиты.
\end{itemize}

% ---------------------------------------------------------------------------
\section*{Упражнения}

\begin{esercizio}[1]\label{es:01-metadati}
Перечислите метаданные, которые интернет-провайдер может наблюдать, когда его клиент читает статью на новостном сайте по HTTPS. Какие данные остаются видимыми, если клиент использует VPN? А если Tor?
\end{esercizio}

\begin{esercizio}[1]\label{es:01-termini}
Отнесите каждый из следующих примеров к анонимности, псевдонимности или ни к одной из этих категорий: аккаунт в социальной сети под вымышленным именем; неподписанное сообщение на публичной доске объявлений; письмо, подписанное инициалами; пожертвование наличными, опущенное в ящик для сбора средств.
\end{esercizio}

\begin{esercizio}[2]\label{es:01-insieme}
За один час сайт посещают 50 разных пользователей Tor. Один из них публикует комментарий с орфографической ошибкой, типичной для носителей итальянского языка, причём комментарий появляется в 3:00 по итальянскому времени. Как это изменяет множество анонимности пользователя? Что этот пример показывает о роли поведения самого пользователя?
\end{esercizio}

\begin{esercizio}[2]\label{es:01-vpn}
Сравните коммерческий VPN и Tor с точки зрения того, кто способен связать пользователя с его пунктом назначения. Где в каждом случае сосредоточено доверие?
\end{esercizio}

% ---------------------------------------------------------------------------
\section*{Дополнительная литература}
\begin{itemize}
  \item A.~Pfitzmann, M.~Hansen, \titolo{A terminology for talking about privacy by data minimization}~\cite{pfitzmann2010terminology}: базовая терминология.
  \item P.~Syverson, \titolo{A Peel of Onion}~\cite{syverson2011peel}: история onion routing, рассказанная одним из его создателей.
  \item Tor Metrics, \url{https://metrics.torproject.org/}: оценки числа пользователей по странам.
\end{itemize}
